\documentclass[11pt]{article}
\usepackage[letterpaper,margin=.9in]{geometry}
\usepackage{amsmath,amssymb,amsthm,hyperref}
\hypersetup{hidelinks,pdftitle={The separated zero pairs in Proposition 4}}
\newcommand{\R}{\mathbb R}
\newcommand{\e}{\mathrm e}
\newcommand{\ii}{\mathrm i}
\newcommand{\wt}{\mathrm d}
\newcommand{\ind}{\mathbf 1}
\newtheorem{proposition}{Proposition}
\title{The separated zero pairs in Proposition 4}
\author{Research notes}
\date{}
\begin{document}
\maketitle
Assume Dirichlet GRH. Write $L=\log p$, $y=cL^2$, $\lambda=\log y$,
$\chi=(\cdot/p)$ and $w(v)=\e^{-2|v|}$. All limits in $p$ hold with
$c,T>0$ fixed unless specified. The zeros $\tfrac12+\ii\gamma$ of
$L(s,\chi)$ are indexed with multiplicities and with both signs of the
ordinate. Proposition 4 of the accompanying note gives constants
$T_c,b_c>0$ for which $n_p>y$ implies
\[
 F_p(y,T_c)=\int_{\R}\left|\sum_{|\gamma|\le T_c}
           \e^{\ii(\lambda+v)\gamma}\right|^2 w(v)\,\wt v\ \ge b_cL^2 .
\]
To obtain $n_p=o(L^2)$ it suffices to exclude such primes eventually
for each $c=2^{-j}$, $j\ge1$. No uniformity in $j$ is required.

\section{How far apart must the zeros be?}
Write $W(d)=4/(4+d^2)$. For $H\ge0$ let $C_H$ and $R_H$ denote the
parts of the following ordered pair sum with $|d|\le H$ and $|d|>H$,
respectively, where $d=\gamma-\gamma'$:
\[
 F_p(y,T)=\sum_{|\gamma|,|\gamma'|\le T}W(d)\cos(\lambda d).
 \tag{1}
\]
Indeed $\int_{\R}w(v)\exp(\ii vd)\,\wt v=W(d)$ and the imaginary
parts cancel on reversing the pair.

\begin{proposition}
Uniformly in $0\le H\le1$ at fixed $T$, the sum of the absolute values
of the terms in $C_H$ is $O_T(L^2(H+1/\log L))$. In particular,
$0\le C_{1/\lambda}\ll_{c,T} L^2/\log L$. For an exceptional prime,
\[
 R_{1/\lambda}(y,T_c)\ge (b_c-O_{c,T_c}(1/\log L))L^2 .
 \tag{2}
\]
There is also a fixed $\delta_c>0$ such that
$R_{\delta_c}(y,T_c)\ge b_cL^2/2$ for every sufficiently large
exceptional prime.
\end{proposition}
\begin{proof}
Equations (3.3)--(3.4) of Carneiro et al.\ \cite{cmqr} give, for any
interval $I\subset[-T-1,T+1]$, with the full endpoint multiplicities,
\[
  \#\{\gamma\in I\}=\frac{L}{2\pi}|I|+O_T(L/\log L).
  \tag{3}
\]
On the positive half-axis subtract their counts at the two endpoints.
The $t\log t$ and remaining gamma-factor terms are bounded on a compact
interval, including their limit at zero. One-sided limits handle the
half-weight convention. At zero, taking the positive endpoint down
to zero bounds half the central multiplicity by $O_T(L/\log L)$.
Reflection gives the negative half, and splitting at zero costs at
most the same error. Thus each of the $O_T(L)$ first entries has
$O_T(LH+L/\log L)$ partners. This includes all ordered pairs of entries
at a repeated ordinate. Use $|W\cos|\le1$ and, for the asserted
nonnegativity, $\cos(\lambda d)\ge\cos 1$ in the indicated band.
For fixed $\delta_c$, choose the implicit constant times $\delta_c$
smaller than $b_c/4$ and then let $p$ grow.
\end{proof}
The same counting argument holds for bands about any prescribed
separations. If $E\subset[-2T,2T]$ is a union of $k$ intervals, the
absolute contribution from $d\in E$ is
$O_T\big(L^2(|E|+k/\log L)\big)$. For each first ordinate use (3)
on the translates of those intervals intersected with $[-T,T]$.
Thus a fixed finite list of spacings, each thickened by $o(1)$, can
also be removed.

An absolute-value comparison for the hard cutoff, with an explicit
leading coefficient, is
\[
 F_p(y,T)\le D(T)L^2+O_T(L^2/\log L),\qquad
 D(T)=\frac{8T\arctan T-4\log(1+T^2)}{4\pi^2}.
 \tag{4}
\]
Partial summation twice in (3) gives the integral of $W(t-s)$
over the square $[-T,T]^2$, divided by $4\pi^2$.

\section{A smooth height cutoff and the signed prime powers}
Take a real even $K\in C^\infty_c([-a,a])$, $a>0$, and put
$h(t)=\int K(s)\e^{\ii ts}\,\wt s$. For $v$ in a fixed compact
interval and all sufficiently large $p$, equation (3.1) of
\cite{cmqr}, applied to $h(t)\e^{\ii(\lambda+v)t}$, gives
\begin{align}
 S_h(v)&:=\sum_\gamma h(\gamma)\e^{\ii(\lambda+v)\gamma}
       =-P_0(v)+O_K(1),\nonumber\\
 P_0(v)&=\sum_{n\ge2}a_n(v)\chi(n),\qquad
 a_n(v)=\frac{\Lambda(n)}{\sqrt n}K\big(\log(n/y)-v\big).
 \tag{5}
\end{align}
Here the conductor term is $\log(p/\pi)K(-\lambda-v)$ and the
reflected prime term contains $K(-\log n-\lambda-v)$. Both vanish
on the stated interval. The gamma-factor integral is bounded by a
constant times $\int |h(t)|\log(2+|t|)\,\wt t$. The test is entire
of exponential type and rapidly decreasing in horizontal strips,
as required in the explicit formula.

For the moment take $K\ge0$ nonzero and positive in $(-a,a)$, and
write $B_K=\int K(s)\e^{s/2}\,\wt s$. Set
$I=[-3a,2a]$, $I_-=[-3a,-2a]$ and
$G=\int_I |S_h(v)|^2 w(v)\,\wt v$. The same modified statistic has
both the following lower and upper bounds. RH for $\zeta$ and
partial summation give, uniformly on $I$,
\[
 M(v):=\sum_n a_n(v)=B_K\sqrt y\,\e^{v/2}+O_K(\log^2 y).
 \tag{6}
\]
For $n_p>y$ and $v\in I_-$ every $n$ in (5) is at most $y\e^{-a}<y$.
Define $A_J=B_K^2\int_J\e^v w(v)\,\wt v$. Equations (5)--(6) give
\[
 A_{I_-}y-O_K(\sqrt y\log^2 y)\ \le G\le\
 A_Iy+O_K(\sqrt y\log^2 y).
 \tag{7}
\]
The upper bound holds for every prime under GRH. For $n_p>y$ the
exact signed sum in (5) throughout $I$ is
\[
  P_0(v)=M(v)-2\sum_{\substack{n>y\\\chi(n)=-1}}a_n(v).
  \tag{8}
\]
The window lies below $p$ for large $p$. It includes prime powers up
to $y\e^{3a}$, and the second sum is present wherever the support
reaches above $y$. A prime power with negative sign has odd exponent.
In this fixed window any proper prime power has base $<y$ for large
$p$, so its sign on an exceptional prime is actually $+1$, also for
powers above $y$. Thus only primes above $y$ can enter the subtraction.
Without exceptionality proper powers contribute $O_a(\|K\|_\infty)$
in absolute value: the unweighted $\Lambda$ mass of proper powers up
to $X$ is $O(\sqrt X)$ by Chebyshev's bound and summation over the
exponents.

More explicitly, define $J_{mn}=\int_I a_m(v)a_n(v)w(v)\,\wt v$.
Then (5) gives
\[
 G=\sum_{m,n}\chi(mn)J_{mn}+O_K(\sqrt y),\qquad
 \sum_n J_{nn}=O_K(\log y).
 \tag{9}
\]
The diagonal estimate follows from $\Lambda(n)^2\le\Lambda(n)\log n$
and Chebyshev in the window $n\asymp_a y$. Equation (9) retains the
cross terms of both signs. On $I_-$ alone the right side reconstructs
$A_{I_-}y+O_K(\sqrt y\log^2 y)$ for an exceptional prime. Cancellation
in the other $v$-windows cannot remove this positive part of the
integral. Beating (7) from above would already require arithmetic
that prevents the primes in that first window from all splitting.

\section{Transferring a prime-side bound to the original hard cutoff}
For any real even $K$ as in (5), define
$F_h=\int_{\R}|S_h(v)|^2w(v)\,\wt v$. Equations (5)--(6), with
$|K|$ in the majorant in (6), imply
\[
 \limsup_p\frac{\sqrt{F_h}}{L}\le
 \sqrt{4c/3}\int |K(s)|\e^{s/2}\,\wt s.
 \tag{10}
\]
Indeed the pointwise limsup is at most $\sqrt c\,\e^{v/2}$
times the integral on the right before taking the weighted norm.
Absolute continuity and bounded variation suffice for (6) with $|K|$.
The bound $\sum_\gamma |h(\gamma)|=O_h(L)$, from the unit-interval
zero count and rapid decay, permits integration of the limsup by
domination. Finally $\int_{\R}\e^vw(v)\,\wt v=4/3$.

Fix an even, nonnegative Schwartz density $P$ of integral one whose
angular Fourier transform $\widehat P(u)=\int P(t)\e^{-\ii ut}\,\wt t$
is smooth and supported in $[-1,1]$. A normalized square of the
Fourier transform of a real even bump supported in $[-1/2,1/2]$
suffices. Put $m_1=\int|t|P(t)\,\wt t$ and $P_B(t)=BP(Bt)$.
For $B\ge1$ fixed, take $h=\ind_{[-T,T]}*P_B$. Its kernel is
\[
 K(s)=\frac{\sin(Ts)}{\pi s}\widehat P(s/B),\qquad
 |K(s)|\le T/\pi,\quad |s|\le B .
\]
For $H_T=\ind_{[-T,T]}$, equation (3) and the unit-interval count
give
\[
 \lim_{p\to\infty}\frac{1}{L}\sum_\gamma|H_T(\gamma)-h(\gamma)|
 =\frac{\|H_T-h\|_1}{2\pi}\le\frac{m_1}{\pi B}.
\]
For the equality first truncate at a fixed height. On any compact
interval the integrand is piecewise continuous and bounded, and the
tails decrease rapidly. The inequality integrates the symmetric
difference of an interval and its translate, of measure at most
twice the displacement. Since $\int w=1$, the same error bounds
$|\sqrt{F_p(y,T)}-\sqrt{F_h}|/L$ in the limsup. In particular at
$T=T_c$ this smoothing has lower bound
$\liminf_{\rm exc}\sqrt{F_h}/L\ge\sqrt{b_c}-m_1/(\pi B)$ along
any infinite exceptional sequence. One may take $B$ sufficiently
large for each $c$ to make this positive.

Equations (4) and (10) now bound the signed sum for the original
cutoff at each prime before taking the limit:
\[
 \limsup_p\frac{\sqrt{F_p(cL^2,T)}}{L}\le
 \min\left\{\sqrt{D(T)},\
 \inf_{B\ge1}\left[\frac{m_1}{\pi B}
       +\frac{4T}{\pi\sqrt 3}\sqrt c\,\e^{B/2}\right]\right\}.
 \tag{11}
\]
Here we used $\int_{-B}^B\e^{s/2}\,\wt s\le2\e^{B/2}$.
For $0<c\le\e^{-2}$ the choice $B=\tfrac12\log(1/c)$ gives
$O(1/\log(1/c)+T c^{1/4})$ for the left side. This estimate concerns
each fixed $c$, with the limit in $p$ taken first.

For comparison one may use $T_c\ge1$ and the constant from the proof
of Proposition 4 in the original note,
$b_c=c(\int g(v)\e^{v/2}\,\wt v)^2/(4\int |g(v)|^2\e^{2|v|}\,\wt v)$.
By Cauchy--Schwarz on $(-1,0)$, $b_c\le c(1-\e^{-3})/12<c/12$.
For $0<c\le1/2$, both entries of the minimum in (11) exceed
$\sqrt{b_c}$: $D(T_c)\ge D(1)>1/24$ and the exponential term alone
exceeds $\sqrt{c/12}$. The decrease in the displayed upper bound
as $c$ decreases at fixed $T$ gives no contradiction at a fixed $j$.

\section{A second arithmetic input at the individual modulus}
One exact constraint on an actual Legendre symbol is its additive
autocorrelation. Keep the nonnegative $K,I,I_-$ from (7) and put
$P_r(v)=\sum_n a_n(v)\chi(n+r)$ for $r\bmod p$. For large $p$ the
indices in the window are distinct modulo $p$. At each $v$ one has
\[
 \sum_{r\bmod p}|P_r(v)|^2
 =p\sum_n a_n(v)^2-\left(\sum_n a_n(v)\right)^2 .
 \tag{12}
\]
For $m=n$ the complete sum of products of characters is $p-1$.
For $m\ne n$, the substitution $r\mapsto(r+m)/(r+n)$ maps
$\mathbb F_p\setminus\{-n\}$ bijectively to
$\mathbb F_p\setminus\{1\}$. Hence the correlation is $-1$. This
proves (12) while retaining the signed off-diagonal contribution.
In particular, if $E_r=\int_I |P_r(v)|^2 w(v)\,\wt v$, then
\[
 \sum_{r\bmod p} E_r\ll_K p\log y.
 \tag{13}
\]
This is at the given $p$ and can be applied to the exact contribution
of the exceptional translates. Indeed if $n_p>y$, then for every
integer $0\le r\le (1-\e^{-a})y$ and $v\in I_-$, all $n+r$ in
the sum are at most $y$ and have sign $+1$. Consequently all these
$\gg_a y$ translates have $E_r\ge (A_{I_-}+o_K(1))y$ by (6).
Moreover $E_0=G+O_K(\sqrt y)$ by (5). From (7) all these $E_r$
are at least a fixed positive multiple of $E_0$ for large $p$.
Thus (13) implies on an exceptional prime
\[
 y^2\ll_K p\log y,\qquad E_0\ll_K p\log y/y.
 \tag{14}
\]
At $y=c(\log p)^2$, $p\log y/y^2\to\infty$. Even after propagating
the original large value to many translates, the complete correlation
estimate permits it. The direct upper bound $E_0=O_K(y)$ is smaller.

\section{Ranges and the missing inequality}
Kandhil, Languasco and Moree \cite{klm} use the positive-ordinate
statistic. Their Theorem 1 assumes $x\le T\le\e^x$ and
$3\le k\le\exp(x/\log x)$. At $x=cL^2$, $k=p$ satisfies this
conductor condition eventually, while fixed $T_c$ falls below $x$.
Their Theorem 2 requires the same height range and
$k\le\sqrt x/(\log x)^3$, which also excludes $p$. Their additional
Hypothesis 1 starts at $x^\varepsilon\le T<x$ for fixed
$\varepsilon>0$ with the conductor range of Theorem 1. It does not
supply a direct estimate at fixed height.

There is an indirect way to use growing heights. For a band-limited
height cutoff $h$ from (5), set $C_K=\int|K(s)|\e^{|s|}\,\wt s$. For
$U\ge2$ and each fixed $N>0$, one has
\[
 \sqrt{F_h}\le 2C_K\sqrt{F^+_\chi(y,U)}
 +O_{K,N}\big(L/\log L+(L+\log(2U))U^{-N}\big).
 \tag{15}
\]
Indeed the positive-ordinate sum truncated at $U$ and weighted by $h$
is the integral against $K(s)$ of translates by $s$ of the unweighted
sum. Translation costs at most $\e^{|s|}$ in the $L^2(w)$ norm.
Add the reflected sum and subtract the duplicated central entries
using (3). The unit-interval count bounds the discarded Schwartz tail.
This argument uses no monotonicity of $F^+$ in its height.

In the other direction, choose the nonnegative $K$ of (7). For each
$c$ there is a fixed $U_c$ such that every sufficiently large
exceptional prime satisfies, simultaneously for all $U\ge U_c$,
\[
  F^+_\chi(y,U)\ \ge\ \frac{A_{I_-}c}{16C_K^2}\,L^2 .
  \tag{16}
\]
Indeed $F_h\ge G$, and $U_c$ can make the normalized tail in (15)
smaller than a quarter of $\sqrt{A_{I_-}c}$ uniformly in $U\ge U_c$.
For sufficiently large $p$, (7) and the central-zero error leave
at least $\tfrac12\sqrt{A_{I_-}c}\,L$ for the leading term in (15).

Applying the proved Theorem 1 at $U=y$ gives only $F_h\ll_{K,c}L^3$
from (15), above the order $L^2$ in (7). Conversely, an individual
bound $F^+_\chi(y,U)=o_c(L^2)$ at even one height $U=U(p)\to\infty$
at this $y$ gives $\sqrt{F_h}/L\to0$ for each fixed $K$. Approximate the
hard cutoff as in (11) and let $B\to\infty$ after $p\to\infty$ to get
$F_p(y,T)=o_{c,T}(L^2)$. In particular a bound $O_c(UL)$ at some
$1\ll U=o(L)$ with $U\to\infty$ would suffice. For instance, their
additional Hypothesis 1 would allow $U=y^{1/4}$ and would give this
transfer. Its use at that height would be an extra assumption beyond
Dirichlet GRH.

For the original cutoff the close pairs cost $O_{c,T}(L^2/\log L)$.
The fixed-height upper bound (11) still allows the lower bound (2).
A sufficient missing estimate for each fixed dyadic $c$ is
\[
 \limsup_{p\to\infty}
 \frac{R_{1/\log(cL^2)}(cL^2,T_c)}{L^2}<b_c .
\]
A bound restricted to putative exceptional primes with a fixed margin
below $b_c$ would suffice as well. The estimates here give no improved
bound for $n_p$.

\begin{thebibliography}{9}
\bibitem{cmqr} E.~Carneiro, M.~Milinovich, E.~Quesada-Herrera and A.~Ramos,
\emph{Fourier optimization, the least quadratic non-residue, and the
least prime in an arithmetic progression},
\href{https://arxiv.org/html/2404.08380v2#S3}{Section 3}.
\bibitem{klm} N.~Kandhil, A.~Languasco and P.~Moree,
\emph{Beyond the Riemann Hypothesis bounds: A pair-correlation approach
to the least prime in arithmetic progression and the smallest quadratic
non-residue}, \href{https://arxiv.org/html/2607.14515v1#S1.SS1}{Section 1.1}.
\end{thebibliography}
\end{document}
